\section{绪论}

\subsection{研究背景}

视觉语言模型的发展为多模态信息处理开辟了新的可能。以 CLIP 为代表的对比式模型将视觉和语言映射到统一语义空间，BLIP-2、LLaVA 和 Qwen-VL 等生成式模型进一步实现了从图像到自然语言的端到端推理。这些模型在干净图像上性能优异，但在实际部署环境中，图像常因各种因素出现质量退化。当视觉输入受损时，模型的表现会显著下降，这是视觉语言模型从实验室走向实际应用必须面对的问题。

脑电图技术为这一问题的缓解提供了新的思路。EEG 能以毫秒级时间分辨率记录视觉刺激引发的神经响应，已有研究表明其中包含可解码的视觉语义信息。当摄像头的``眼睛''看不清时，人脑这个``生物传感器''能否提供额外的信息？这正是本课程设计希望探索的问题。

本设计将 EEG 定位为视觉语义增强信号——不是在脑电中直接``读''出语言，而是让 EEG 携带的关于图像内容的神经表征在视觉退化条件下补充和修正机器的视觉语义理解。这一研究定位既考虑了 EEG 数据的特点（小样本、低信噪比），也符合当前脑机接口与人工智能交叉领域的实际技术水平。

\subsection{课程设计目标与意义}

本课程设计的目标是：在 Thought2Text 数据集上构建 EEG 辅助的视觉语言理解系统，通过系统性的对照实验验证真实配对 EEG 在六种视觉退化条件下对语义识别和图像描述生成的辅助效果。

从课程实践的角度，本课题综合运用了深度学习模型设计、预训练模型的适配与使用、多模态融合、大模型参数高效微调等技术，是一个较为完整的多模态 AI 系统实践项目。从科学验证的角度，通过 real EEG、shuffled EEG、random EEG 和 vision-only 的多重对照设计，力图为 EEG 是否携带可机器利用的视觉语义信息这一开放问题提供可验证的经验证据。

\subsection{研究内容与技术路线}

本设计完成了以下四项核心工作：

\begin{enumerate}
  \item 搭建了包含六种退化条件和四种 EEG 对照模式的实验评估框架，采用 image-level split 防止训练-测试信号泄漏；
  \item 设计并实现了 A2 temporal-spectral-spatial EEG 语义融合模型，从 EEG 的时间、频谱和空间通道三个维度独立提取特征后融合，通过门控残差机制动态调节 EEG 对退化图像的语义补充强度；
  \item 实现了 VTF token-level 融合方法，将 EEG token 通过 cross-attention 注入 CLIP ViT visual token 序列，使 EEG 信息进入原生 VLM 的 token 处理流程；
  \item 构建了基于 Qwen2-VL 的生成式图像描述模块，使用 LoRA 进行高效微调，对比了五种生成方案，并通过 best-of-N reranking 提升输出质量。
\end{enumerate}

技术路线的设计遵循``先验证后扩展''的原则：利用 CLIP 语义空间作为跨模态桥梁降低对数据量的依赖；先在受约束的分类任务上严格验证 EEG 语义增强的有效性；确认有效后再逐步推进到 token-level 融合和自由文本生成。这种分层验证的策略使每一步都有明确的中间结果可供检查和调试。

\subsection{报告结构}

本报告共十章。第 1 章绪论；第 2 章相关技术基础；第 3 章数据集与任务；第 4 章系统总体设计；第 5 章 A2 EEG 语义融合模型；第 6 章 token-level EEG 视觉增强；第 7 章生成式 VLM 设计；第 8 章实验结果与分析；第 9 章工程实现与代码；第 10 章总结与展望。
